\section{Intervention and Participant Discontinuation/Withdrawal}

This section defines the prospectively specified rules for discontinuing the study
intervention (the device, the drug, or both), the procedures governing a
participant's withdrawal from the study, and the handling of participants lost to
follow-up. Discontinuation of an intervention is distinct from withdrawal from the
study: a participant who stops one or both components of the assigned regimen for a
protocol-defined reason remains in the study, and remains in the randomized
intention-to-treat (ITT) population, for safety and survival follow-up unless
consent for that follow-up is withdrawn. Because this is a randomized, controlled
trial, the rules below apply within the arm to which the participant was
randomized; the device stopping rules apply to Arm A, and the drug stopping rules
are specified separately for daraxonrasib in Arm A and for modified FOLFIRINOX in
Arm B.

\subsection{Discontinuation of Study Intervention}

\paragraph{Device stopping rules.}
For a participant randomized to Arm A, the investigational device is discontinued
for the index procedure, with conversion to the appropriate fallback technique,
whenever any of the following prespecified conditions occurs. These rules implement
the Physical AI stopping criteria of 21 CFR $\S$312.23(g)(i)(v) and the
divergence-reporting threshold of $\S$312.32(g) \cite{kawchak2026adapt312}, and the
federated hash-chained audit trail is preserved across each event so that the
behavior can be reconstructed across the eight-site fleet.

\begin{itemize}[leftmargin=1.4em,itemsep=1pt,topsep=2pt]
\item \textit{Conversion to open or manual technique.} Any intra-operative judgment
by the supervising surgeon that patient safety or oncologic adequacy is better
served by manual robotic or open conversion (including uncontrolled bleeding,
unfavorable anatomy, or loss of a critical view) ends device-directed operation
immediately; conversion is always available and is never itself a protocol
violation.
\item \textit{Repeated force-limit excursions.} Repeated excursions against the
$\leq 3$ N per-arm or $\leq 18$ N cumulative cross-arm tip-force caps, or any
sustained or consecutive pattern of force-clamp activations indicating that the
case cannot proceed within the force envelope, trigger discontinuation.
\item \textit{No-fly-zone violations.} Any breach of a vascular no-fly zone around a
named vessel (superior mesenteric vein or portal vein at 1.0 mm; hepatic artery,
celiac axis, or superior mesenteric artery at 1.5 mm), or a recurring pattern of
no-fly approaches, ends device-directed dissection in that field.
\item \textit{Sim-to-real divergence beyond twice tolerance.} Any instance in which
the observed device behavior diverges from the validated simulation prediction by
more than twice the sim-to-real gap tolerance (a trajectory deviation greater than
3 mm or a force discrepancy greater than 0.8 N) triggers discontinuation and a
report under $\S$312.32(g) \cite{kawchak2026adapt312}.
\item \textit{Hard-stop E-stop event.} Any hard-stop cross-arm E-stop (the
$\leq 3$ ms full-platform halt, whether software- or hardware-initiated) ends the
device-directed segment; resumption with the device requires that the cause be
identified and that the system re-pass the pre-procedure safety matrix
($\S$6.2), failing which the case is completed by conversion.
\end{itemize}

Following any device discontinuation, the operation is completed by the most
appropriate conventional means, the federated hash-chained audit trail is preserved
across the event window ($\S$10.10), and the event is reviewed by the Physical AI
Safety Review Committee and the data and safety monitoring board (DSMB) before the
device is used in a subsequent case at any site. A participant whose index procedure
is converted is analyzed in Arm A as randomized ($\S$9).

\paragraph{Drug stopping rules.}
For a participant in Arm A, daraxonrasib is discontinued for an unacceptable
toxicity that does not resolve with protocol dose interruption and reduction, for a
confirmed International Study Group of Pancreatic Surgery (ISGPS) Grade C
postoperative pancreatic fistula \cite{Bassi2017} (which precludes restart), for
pregnancy, for intercurrent illness that contraindicates continuation, or at the
participant's request. For a participant in Arm B, modified FOLFIRINOX is managed,
interrupted, dose-reduced, or discontinued according to standard toxicity
management for that regimen \cite{Conroy2018}. In either arm, dose interruptions and
reductions for lower-grade toxicity follow the applicable algorithm and do not by
themselves mandate permanent discontinuation. A participant who permanently
discontinues the assigned drug, in either arm, remains in the study and in the ITT
population for survival follow-up to the 24-month endpoint.

\subsection{Participant Discontinuation/Withdrawal from the Study}

A participant may withdraw from the study at any time, for any reason, without
penalty and without loss of any benefit to which the participant is otherwise
entitled, and a participant may separately withdraw consent for continued
study-specific procedures while permitting ongoing survival follow-up through the
medical record. The investigator may withdraw a participant for safety, for
intercurrent illness, for significant non-compliance that compromises participant
safety or data integrity, for loss of eligibility, or at the sponsor's direction if
the study is terminated. The reason for and date of any discontinuation or
withdrawal are recorded in the case report form (CRF), and a participant who
withdraws is asked to complete a final safety assessment where feasible.

Data and specimens collected up to the point of withdrawal are retained and analyzed
in accordance with the consent and applicable regulation; withdrawal of consent for
future participation does not require deletion of data already collected. The
relationship of discontinuation and withdrawal to the analysis populations (the
ITT, modified ITT, per-protocol, and safety populations) is defined in the
statistical considerations ($\S$9). Because this is a randomized ITT trial,
randomized participants are analyzed in the arm to which they were randomized,
regardless of the intervention actually received or completed, and are not replaced;
no participant who discontinues an intervention or withdraws is substituted by a
newly randomized participant.

\subsection{Lost to Follow-Up}

A participant is considered lost to follow-up only after the site has made and
documented at least three contact attempts by telephone across separate days,
followed by a written contact attempt to the last known address, over the 24-month
overall-survival window. Each attempt is recorded with its date and outcome in the
CRF. Before declaring a participant lost to follow-up, the site attempts, with
appropriate authorization, to ascertain vital status from the medical record, the
referring provider, and permissible public sources, because overall survival is a
key endpoint of the study. A participant who cannot be reached after these attempts
is classified as lost to follow-up, with the last date known alive recorded. The
handling of the resulting missing data, including the censoring rules for
progression-free and overall survival and the sensitivity analyses used to assess
the impact of missingness, is specified in $\S$9.
